\section{Common-policy certification without a strict action gap}\label{sec:central45}
Set
\begin{align}
 D_C&=K\sum_{s=0}^{T-1}(s+1)\beta^s,\label{eq:DC45}\\
 D_J&=\max_{1\le h\le T}\left\{2R\sum_{s=0}^{h-1}(s+1)\beta^s+2Gh\beta^h\right\}.\label{eq:DJ45}
\end{align}
If two policies differ by at most $\theta$ in action-distribution total variation at every state and date, their operating values differ by at most $D_J\theta$, and their costs by at most $D_C\theta$. To see this, subtract their Bellman equations. A cost at distance $s$ can be affected by $s+1$ action changes, and a terminal operating payoff by $h$ changes. Using the range, rather than the absolute magnitude, of the action values gives the stated constants. Total variation is one half of the $\ell_1$ distance.

\begin{theorem}[Backward repair]\label{thm:repair45}
For $p\in\cF(\varepsilon+\delta)$, $\delta\ge0$, there exists a Borel common policy $\mathcal R_\varepsilon p\in\cF(\varepsilon)$ such that $J_t^{\mathcal R_\varepsilon p}\ge J_t^p$ at every state and date, and
\begin{equation}\label{eq:rho45}
 \sup_{t,x}\TV((\mathcal R_\varepsilon p)_t,p_t)\le
 \rho_\varepsilon(\delta):=\frac{\delta}{(1-\beta)\varepsilon+\delta}.
\end{equation}
Its cost exceeds that of $p$ by at most $D_C\rho_\varepsilon(\delta)$. In particular,
\begin{equation}\label{eq:mod45}
 0\le\KR(\nu,\varepsilon)-\KR(\nu,\varepsilon+\delta)
 \le D_C\rho_\varepsilon(\delta).
\end{equation}
\end{theorem}
\begin{proof}
Fix the repaired continuation backward. At $(t,x)$ put $q_a=r_t(x,a)+\beta P_t^aJ_{t+1}^{\mathcal R_\varepsilon p}(x)$, $q_p=\sum_a p_t(a\mid x)q_a$, $q_*=\max_aq_a$, and $b=V_t(x)-\varepsilon$. The induction hypothesis gives $q_p\ge V_t-\varepsilon-\delta$ and $q_*\ge V_t-\beta\varepsilon$. If $q_p\ge b$, retain the original action distribution. Otherwise mix it with a first maximizing action in proportion
\[
 \alpha=\frac{b-q_p}{q_*-q_p}.
\]
The denominator is at least $(1-\beta)\varepsilon+(b-q_p)$, so $\alpha\le\rho_\varepsilon(\delta)$. The repaired value is $b$, weakly above both $q_p$ and the original value. Every operation is Borel, and the division is used only with a positive denominator. This constructs one continuation, not a mixture of persistent plans. The perturbation bound gives the cost inequality; apply it to approximate minimizers to obtain \eqref{eq:mod45}.
\end{proof}
The same formula defines $\mathcal R_\varepsilon q$ for an arbitrary Borel proposal $q$, because the repaired continuation always supplies $q_*\ge V_t-\beta\varepsilon$. A small uniform cost bound requires a small violation, but existence of the repair does not. This distinction will be used for continuous states.

\subsection{Tie-safe support envelopes}\label{sec:support45}
For a fixed real price $\lambda$, define $W_{T,\lambda}^-=W_{T,\lambda}^+=0$ and
\begin{equation}\label{eq:W45}
 W_{t,\lambda}^{-}=\min_a\{h_t-\lambda d_t+\beta P_t^aW_{t+1,\lambda}^{-}\},\qquad
 W_{t,\lambda}^{+}=\max_a\{h_t-\lambda d_t+\beta P_t^aW_{t+1,\lambda}^{+}\}.
\end{equation}
\begin{theorem}[Bellman supports at and near ties]\label{thm:support45}
For every Borel common policy, without a uniqueness or positive-gap assumption,
\begin{equation}\label{eq:support45}
 W_{t,\lambda}^{-}\le S_t^p-\lambda D_t^p\le W_{t,\lambda}^{+}.
\end{equation}
Consequently any feasible policy satisfies all inequalities \eqref{eq:support45} together with $0\le D_t^p\le\varepsilon$. For $\lambda\ge0$, let
\begin{equation}\label{eq:Z45}
 Z_{T,\lambda}=\lambda g,\qquad
 Z_{t,\lambda}=\max_a\{\lambda r_t-k_t+\beta P_t^aZ_{t+1,\lambda}\}.
\end{equation}
Then $\ell_{t,\lambda}:=\lambda(V_t-\varepsilon)-Z_{t,\lambda}$ is a lower bound for $C_t^p$, and it is independent of $H$. For any finite set $\Lambda\subset[0,\infty)$,
\begin{equation}\label{eq:supportlower45}
 \int\max\{0,\max_{\lambda\in\Lambda}\ell_{0,\lambda}(x)\}\,\nu(dx)
 \le \KR(\nu,\varepsilon).
\end{equation}
\end{theorem}
\begin{proof}
Subtract $\lambda$ times \eqref{eq:D45} from \eqref{eq:S45}. Backward induction and the fact that a convex combination lies between its smallest and largest action value give \eqref{eq:support45}. Direct substitution yields $W_{t,\lambda}^+=H_t-\lambda V_t+Z_{t,\lambda}$. Thus $C_t^p\ge H_t-W_{t,\lambda}^+-\lambda\varepsilon=\ell_{t,\lambda}$. Pointwise maxima and nonnegativity preserve the lower inequality before integration. No transition has been approximated in this argument.
\end{proof}
For a finite model, each inequality in \eqref{eq:support45} is linear in $D,S$. The frozen enhanced formulation uses $\lambda=0$ and $\lambda=\pm2^j$, $j\in\{-2,0,2,4,6,8,10,12\}$. The prices are fixed before optimization. Their finiteness makes them inexpensive cuts, not a claim of exact global duality.

\begin{proposition}[Affine cones for approximate ties]\label{prop:affine45}
Suppose finite constants $\lambda_-\le\lambda_+$ and nonnegative numbers $e_t$ satisfy
\[
 \lambda_-d_t(x,a)-e_t\le h_t(x,a)\le\lambda_+d_t(x,a)+e_t
 \quad\hbox{for every }(x,a).
\]
With $B_T=0$ and $B_t=e_t+\beta B_{t+1}$,
\begin{equation}\label{eq:affine45}
 \lambda_-D_t^p-B_t\le S_t^p\le\lambda_+D_t^p+B_t
\end{equation}
for every policy. In particular the result permits $d_t=0$ and nonzero $h_t$. If, instead, only actions with $d_t\ge\tau>0$ are classified as separated, every feasible policy assigns their total probability at most $\varepsilon/\tau$.
\end{proposition}
\begin{proof}
Subtract $\lambda_+D$ from $S$ in the joint Bellman recursion and bound the stage term by $e_t$. Induction gives the upper inequality. The lower inequality is identical with $-e_t$. For the last claim, nonnegativity in \eqref{eq:D45} gives $\sum_a p_t(a\mid x)d_t(x,a)\le\varepsilon$.
\end{proof}
The affine intercept records the cost of a near-tie group instead of hiding it inside an inverse minimum gap. It need not be small. In particular, an exact operating tie may allow a substantial implementation-cost difference. The frozen solver uses the Bellman supports, not this newly stated affine-cone variant; the latter is an analytic strengthening available for subsequent implementations.

\subsection{The strict-gap specialization}\label{sec:quadratic45}
The earlier regret-scaled rate remains a useful specialization. With an operating-optimal reference and strictly positive nonreference disadvantages, put
\[
 \begin{gathered}
 \underline d=\min_{t,i,a\ne a^*}d_{tia},\qquad \chi=\min\{1,\varepsilon/\underline d\},\\
 \lambda_- =\min\{0,\min h_{tia}/d_{tia}\},\qquad
 \lambda_+ =\max\{0,\max h_{tia}/d_{tia}\}.
 \end{gathered}
\]
Feasible nonreference mass is at most $\chi$, and $\lambda_-D\le S\le\lambda_+D$. With $\Delta P_{tia}=P_i^a-P_i^{a^*}$, define $e_{D,T-1}=0$ and, for $t<T-1$,
\begin{align}
 e_{D,t}&=\frac{\beta\varepsilon^2}{4}\max_i\sum_{a\ne a^*}
             \frac{\|\Delta P_{tia}\|_1}{d_{tia}},&
 E_{D,t}&=e_{D,t}+\beta E_{D,t+1},\label{eq:ED45}\\
 E_{S,t}&=(\lambda_+-\lambda_-)e_{D,t}+\beta E_{S,t+1},&
 \overline E_D&=\max_tE_{D,t},\nonumber
\end{align}
with zero terminal errors. An exact optimizing root-LP solution, repaired toward the operating reference, has width bounded by
\begin{equation}\label{eq:quadratic45}
 E_{S,0}+D_C\chi\frac{\overline E_D}{(1-\beta)\varepsilon+\overline E_D}.
\end{equation}
The bound follows from the McCormick product error, at most one quarter of the product of the two interval widths, followed by the rare-departure repair. The full derivation is retained in the historical technical material. At fixed primitives, fixed $\beta<1$, and $\underline d>0$, \eqref{eq:quadratic45} is $O(\varepsilon^2)$. A numerical LP optimality loss and directed endpoint errors must be added when an implemented certificate is compared to that exact-LP statement. A solver status code does not set those losses to zero.

The rate neither covers ties nor remains uniform when $\underline d$ vanishes along state refinements. Formula \eqref{eq:affine45} and Theorem~\ref{thm:search45} below give valid alternatives without declaring the strictly separated theorem uniform. The tables report the actual constants and exclude tied cases from the inverse-gap formula.

\subsection{A complete multi-action search and its numerical floor}\label{sec:tree45}
For finite $X$, use one coordinate for every nonreference probability at every state and date. A node is a box in these $d=nT(m-1)$ coordinates, intersected with the local probability simplexes. The reference probability is the remaining mass. The lower corner belongs to every nonempty such row because the sum of its nonreference lower bounds cannot exceed one.

Restricted Bellman minimization of regret and implementation cost over these local simplexes gives, respectively, a necessary operating test and a cost lower bound. These restricted dynamic programs are not the final constrained optimization. They are admissible relaxations of all the common policies in the box. The linear formulation additionally imposes \eqref{eq:D45}--\eqref{eq:S45}, the support cuts, and product envelopes. In the disaggregated formulation the products also obey
\begin{equation}\label{eq:simplexRLT45}
 \sum_{a\in A} w_{tia,j}=z_{t+1,j},\qquad w_{tia,j}=p_{tia}z_{t+1,j}
 \quad\hbox{before relaxation},
\end{equation}
for each retained continuation coordinate $z=D,S$. The aggregate formulation instead relaxes $p_{tia}(\Delta P_{tia}z_{t+1})$. These are different relaxations of the same Bellman equations. Disaggregation costs more variables and arithmetic per node.

A floating-point LP supplies only a candidate policy and multipliers. For a minimization LP $\min c'x$ subject to $Ax=b$, $Gx\le h$, and $l\le x\le u$, any equality multipliers $y$ and nonnegative inequality multipliers $\mu$ give
\begin{equation}\label{eq:dual45}
 b'y-h'\mu+\sum_j\min\{(c-A'y+G'\mu)_j l_j,
                              (c-A'y+G'\mu)_j u_j\}
\end{equation}
as a lower bound. Directed arithmetic encloses each term outward. Zero multipliers remain valid if a proposal is unavailable. A floating-point infeasibility flag alone never removes a box.

The incumbent is the smallest exactly evaluated repaired-policy cost. Every split retains both children. Infeasible leaves require a rational necessary inequality; all other terminal leaves retain their lower endpoint. The global lower endpoint is the minimum of those leaf endpoints and the incumbent. This full cover is part of the proof object.

\begin{theorem}[Gap-free finite stopping with explicit precision]\label{thm:search45}
Consider a surviving probability box of maximum nonreference coordinate width $w$. Suppose each restricted Bellman scalar result is rounded outward with error at most $q=2^{-b}$ per date, and all other node removals and additional lower bounds are valid. Put
\begin{equation}\label{eq:H45}
\begin{gathered}
 \theta(w)=\min\{1,(m-1)w\},\qquad \gamma_b=2^{-b}\Ssum T,\\
 \mathcal H(w,b)=D_C\theta(w)+\gamma_b+
                   D_C\rho_\varepsilon(D_J\theta(w)+\gamma_b).
 \end{gathered}
\end{equation}
The cost of exact backward repair of the lower-corner policy minus the rounded restricted cost lower bound is at most $\mathcal H(w,b)$. For every $\eta>0$, choose $b$ and $w_*>0$ with $\mathcal H(w_*,b)\le\eta$. Best-bound bisection, with a widest-coordinate split at least once every eight levels of a path, then terminates at gap $\eta$ without node or time caps. One sufficient full-tree depth bound is
\begin{equation}\label{eq:depth45}
 8d\left\lceil\log_2(1/w_*)\right\rceil,
\end{equation}
where the ceiling is replaced by zero if $w_*\ge1$. No minimum positive action disadvantage enters these statements.
\end{theorem}
\begin{proof}
The detailed argument is in Appendix~\ref{app:search45}. Restricted Bellman minimization attains minimum regret simultaneously at every state with one Markov policy in the box. Its exact minimum regret is at most $\varepsilon+\gamma_b$ in a surviving node. All policies in that box differ by at most $\theta(w)$ in action-distribution total variation. The corner therefore violates operating feasibility by at most $D_J\theta(w)+\gamma_b$. Its cost is at most $D_C\theta(w)+\gamma_b$ above the rounded restricted minimum. Repair proves \eqref{eq:H45}. A small enough box cannot require further splitting when its repaired policy has been considered. Widest-coordinate splits force all widths below $w_*$ by the stated depth, independently of heuristic intervening splits. There are finitely many nodes up to that depth.
\end{proof}
The theorem is uniform over finite models with common bounds $R,K,G,T,m,\beta,\varepsilon$ for the stopping modulus, even when action gaps vanish. It is not uniform in computational dimension: the depth contains $d$, and the number of tree nodes is exponential in depth. Nor does it give a joint approximation theorem merely by substituting a finer state grid for the original kernel.

The frozen code uses two grids: $2^{-44}$ for node-bound arithmetic and $2^{-40}$ for deployed nonreference probabilities. It also uses an operating-reference repair with a rounding guard, rather than the exact first-maximizer repair in Theorem~\ref{thm:repair45}. The following additional bound charges this implementation difference.

\begin{corollary}[Quantized operating-reference repair]\label{cor:quantized45}
Let $b_v$ denote the bound precision and $b_p$ the deployed-probability precision. Define
\begin{equation}\label{eq:quantities45}
 \begin{gathered}
 s=(1-\beta)\varepsilon,\quad B=\max_{t,i,a}d_{tia}+\beta\varepsilon,
 \quad \kappa=(m-1)2^{-b_p},\quad \mu=\kappa B,\\
 \delta=D_J\theta(w)+\gamma_{b_v},\qquad
 \Delta=\delta+\mu\Ssum T.
 \end{gathered}
\end{equation}
Assume $\mu<s$. At each row, mix the proposal toward the operating reference until its regret is at most $\varepsilon-\kappa\max_a|q_a-q_{a^*}|$, where $q_a=d_{tia}+\beta P_i^a\widehat D_{t+1}$ uses the already repaired continuation. Round each nonreference probability downward to a multiple of $2^{-b_p}$ and assign the remaining mass to the reference. This is the frozen repair rule. Its repaired lower-corner cost minus the rounded node cost lower bound is at most
\begin{equation}\label{eq:HQ45}
 \mathcal H_Q(w,b_v,b_p)=D_C\theta(w)+\gamma_{b_v}
       +D_C\min\left\{1,\kappa+\frac{\Delta}{s-\mu+\Delta}\right\}.
\end{equation}
The finite-depth conclusion of Theorem~\ref{thm:search45} holds with $\mathcal H_Q$ in place of $\mathcal H$. Arbitrary positive target accuracy requires choosing both precisions sufficiently large; fixing either grid need not give a vanishing sufficient floor.
\end{corollary}
\begin{proof}
Appendix~\ref{app:search45} bounds the accumulated regret increase due to rounding by $\mu\Ssum T$, the mixing fraction by $\Delta/(s-\mu+\Delta)$, and the extra row total variation by $\kappa$. Apply the implementation-cost modulus to the repaired lower corner. The reference slack $s$ is positive without any strict action gap.
\end{proof}
The frozen pair $(b_v,b_p)=(44,40)$ gives sufficient floors below $10^{-3}$ on all sixteen recorded models; the largest is about $3.715\times10^{-4}$ for Q3. The corresponding exact-repair floor is smaller and is not substituted for this implemented-repair bound. The supplementary tables give both versions and their sufficient widths. Actual certificate endpoints remain valid independently of these conservative bounds. The pre-existing fixed-precision sources and measurements are unchanged; no adaptive-precision experiment is claimed.

At $\varepsilon=0$, a different exact argument applies. Because all terms in \eqref{eq:D45} are nonnegative, every feasible policy uses only operating-maximizing actions at every state and date. Conversely those actions give zero regret by backward induction. Minimizing implementation cost by Bellman recursion within that action set therefore solves the zero-allowance problem, including ties. This statement does not require taking the singular limit of \eqref{eq:rho45}. The retained finite-horizon allowance-ladder extension likewise has its own assumptions when $\beta=1$.
